\documentclass[tikz,border=10pt]{standalone}
\usetikzlibrary{calc}

\begin{document}
\begin{tikzpicture}
    % Define coordinates
    % B is the origin (0,0)
    \coordinate (B) at (0,0);
    % A is directly above B
    \coordinate (A) at (0,3);
    % E is on the horizontal line from B
    \coordinate (E) at (1.2,0);
    
    % F is the reflection of B across the line AE.
    % This corresponds to folding the triangle ABE along AE.
    % First, find the projection of B onto AE
    \coordinate (P) at ($(A)!(B)!(E)$);
    % Then F is such that P is the midpoint of BF
    \coordinate (F) at ($(B)!2!(P)$);
    
    % G is the intersection of the line extending AF and the line BC.
    % We define a point X far to the right to define the line BC.
    \coordinate (X) at (10,0);
    \coordinate (G) at (intersection of A--F and B--X);
    
    % C is further to the right on the bottom line
    \coordinate (C) at (7.5,0);
    % D is the top-right corner, AD is horizontal
    \coordinate (D) at (4,3);
    
    % Draw dashed lines
    \draw[dashed] (A) -- (B);
    \draw[dashed] (B) -- (E);
    \draw[dashed] (F) -- (G);
    
    % Draw solid lines
    \draw (A) -- (E);
    \draw (A) -- (F);
    \draw (E) -- (F);
    \draw (A) -- (D);
    \draw (D) -- (C);
    \draw (E) -- (C);
    
    % Draw right-angle mark at B
    \draw (0, 0.3) -- (0.3, 0.3) -- (0.3, 0);
    
    % Label the points
    \node[left] at (A) {$A$};
    \node[below left] at (B) {$B$};
    \node[below] at (E) {$E$};
    \node[below] at (G) {$G$};
    \node[below right] at (C) {$C$};
    \node[above] at (D) {$D$};
    \node[above right] at (F) {$F$};
\end{tikzpicture}
\end{document}
